\section{Conclusion}\label{sect:Conclusion}


In this theoretical investigation we have developed a new approach for considering the air layer trapped between an impacting drop and a liquid pool, which permits the study of rebound off a deep liquid bath in a more physically complete manner relative to methodologies which ignore the presence of the gas altogether. Through evolving each of the three regions separately and providing a method of coupling through the pressure obtained from a thin film approximation for the air, we have shown that it is possible to obtain good agreement with results obtained via direct numerical simulation.

The lubrication-mediated model that we have introduced has the capability of capturing the height and pressure within the air layer through the solution of an elliptic thin film problem over a dynamic boundary, which is novel in this context. We are able to solve for solid impactors and for deformable droplets in the small deformation limit. Our model is also restricted to small bath deformations and small viscous effects.

We have shown through an in-depth comparison for a single rebound case that the lubrication-mediated model is able to qualitatively match with the DNS system, obtaining similar behaviours and dynamics. A highlight is the comparison of air layer thicknesses shown in the previous section, which has good agreement with the morphology of air layers seen in droplet rebound experimental \cite{tang2019bouncing} and DNS studies. 

 Across a range of parameters we have uncovered that the lubrication-mediated model is able to capture behaviour of models based on the kinematic match approach, matching qualitatively with the respective 3D reduced methods \cite{galeano2021capillary, alventosa2023inertio}, yet providing additional insight into the air layer dynamics without increased computational effort. 
 We have discussed the numerical stiffness observed in our system, and have proposed viscous damping methods to mitigate this issue. Further investigation is needed to improve the method and adapt it to regimes which permit larger deformation of the drop and liquid bath. 

 There are several avenues to extend this work, and the related 3D rigid-sphere axisymmetric computations in \cite{phillips2024lubrication}. Two direct extensions are three-dimensional droplet impacts and two-dimensional non-symmetric impacts.

 Finally, we have demonstrated the robustness of the lubrication-mediated model, and its applicability to repeated bouncing behaviour under sustained forcing as first discovered  by Couder \textit{et al} \cite{couder2005walking}. Future avenues of work  would be to investigate the results of the lubrication mediated model in Faraday pilot wave studies such as for walkers where almost vertical impacts capture the dynamics \cite{galeano2017non}, and in dynamics where the vertical motion is known to play an important role such as in corrals \cite{durey2020faraday}, tunnelling \cite{nachbin2017tunneling} and non-resonant behaviour \cite{primkulov2025nonresonant}. 

